\section{Discussion}
\label{sec:discussion}

\subsection{Implications for Practitioners}

Our findings provide actionable guidance for LLM selection in software engineering:

\textbf{Speed-Critical Workflows:} For rapid prototyping or real-time assistance, GPT-5.1 offers the best combination of speed (44s avg) and quality (perfect score). GPT-4o provides nearly equal quality at 25\% faster average times.

\textbf{Cost-Conscious Deployment:} Gemini-2.5 Flash achieves 80\% scores at the fastest speeds (51s avg), making it cost-effective for high-volume applications where perfect accuracy is not required.

\textbf{Research Tasks:} Models showed most variation in research synthesis. GPT-5.1 and Deepseek-Chat achieved excellent results, while others struggled with citations. For research-heavy workflows, these models are recommended.

\textbf{Tool Budgeting:} Our finding of no linear correlation between tool usage and score ($r=0.077$; see Section~\ref{sec:threats}) suggests that API costs from excessive tool calls may not improve outcomes. We recommend:
\begin{itemize}
    \item Implementing tool call limits (e.g., 50 calls max) with early termination
    \item Adding loop detection to identify repetitive sequences
    \item Monitoring cost-per-task metrics alongside accuracy
\end{itemize}

\textbf{Budget Implications:} The 53$\times$ cost variance means a \$100/month workload on Gemini-2.5 Flash would cost \$5,300/month on GPT-5.2. For high-volume applications, model selection has significant budget implications.

\subsection{Model Selection Framework}
\label{sec:selection}

Synthesising the performance, efficiency, and cost results, we distil a constraint-driven selection framework (Table~\ref{tab:selection}). Rather than a single ``best'' model, the appropriate choice depends on the dominant deployment constraint, because models with identical quality differ by up to 22--53$\times$ in time and cost.

\begin{table}[htbp]
\centering
\small
\caption{Constraint-driven model selection framework derived from our results.}
\label{tab:selection}
\begin{tabular}{@{}lll@{}}
\toprule
\textbf{Dominant constraint} & \textbf{Recommended} & \textbf{Rationale} \\
\midrule
Speed-critical & GPT-5.1; Gemini-2.5 Flash & fastest at top quality \\
Quality-critical & GPT-5.1, Gemini-3 Pro, & perfect 10/10 score \\
 & Deepseek-Chat, GLM-4.7 & \\
Cost / high-volume & Gemini-2.5 Flash; Deepseek & lowest cost tier \\
On-premise / privacy & GLM-4.7 & best open-weight \\
Research synthesis & GPT-5.1; Deepseek-Chat & reliable citations \\
\bottomrule
\end{tabular}
\end{table}

\subsection{Surprising Findings}

\textbf{Extreme Efficiency Variance:} The most striking finding is the 49$\times$ variance in tool efficiency among models achieving identical scores. Gemini-3 Flash used 917 tool calls for bug-fix (vs GPT-5.1's 3 calls) yet both achieved EXCELLENT. This represents a potential 100$\times$ cost difference for equivalent outcomes.

\textbf{Two Inefficiency Mechanisms:} We identified distinct failure modes:
\begin{itemize}
    \item \textit{Loop inefficiency}: Agents repeating tool sequences without progress (Gemini-3 Flash pattern)
    \item \textit{Inference inefficiency}: Correct solutions generated slowly (Qwen3-VL pattern)
\end{itemize}
These require different mitigation strategies: loop detection for Type A, and model selection/caching for Type B.

\textbf{Framework-Model Incompatibility:} Kimi-K2.5's failure was caused by malformed parallel tool calls, not capability limitations. This highlights the importance of testing framework compatibility before deployment.

\textbf{Perfect Coding Scores:} All 11 models achieved 100\% success on coding tasks (bug-fix, feature, refactor). This indicates that SE coding tasks may be approaching saturation for current LLMs, suggesting future benchmarks need increased difficulty.

\textbf{Research Challenges:} The research task showed the lowest success rate (90.9\%) and highest variance. Citation management and source synthesis remain challenging for LLMs.

\subsection{Threats to Validity}
\label{sec:threats}

We organise threats following the standard internal, external, construct, and conclusion validity categories.

\textbf{Internal Validity.} Each model--task combination was executed once ($N=1$); a single run cannot capture the stochastic variance of LLM generation, so individual durations and tool-call counts are point observations rather than expected values. We also did not control decoding parameters (e.g., temperature) across providers---several models expose them differently or not at all---and task difficulty is not guaranteed to be uniform across the five categories. These factors limit causal attribution of observed differences to the models alone.

\textbf{External Validity.} Our five tasks are synthetic and controlled; while designed to be representative of common developer activities, they omit the ambiguity, scale, and evolving context of real-world software engineering. All tasks are Python-centric and English-only, and the DevOps-oriented prompts (Zrb-Flow, Docker, Kubernetes) may not generalise to other domains or programming languages. We did not evaluate Anthropic's Claude models \cite{claude3} due to budget constraints, leaving one major provider uncovered. Because LLM capabilities evolve rapidly, results reflect model versions available as of February 2026 and may not hold for later releases.

\textbf{Construct Validity.} Automated, deterministic verification yields objective and reproducible grades but cannot assess qualitative attributes such as code readability, maintainability, or architectural elegance. Our cost metric (Eq.~\ref{eq:cost}) assumes a constant token footprint per tool call; actual billing depends on prompt and context length, so reported costs are relative rather than absolute. The absence of a human baseline means we cannot position model performance relative to expert developers.

\textbf{Conclusion Validity.} With $n=54$ model--task observations and near-ceiling success rates, our correlation estimates carry wide uncertainty. The tool-usage/score Pearson correlation of $r=0.077$ has a 95\% confidence interval of $[-0.20, 0.34]$, which spans zero; we therefore interpret the absence of a tool--success relationship as a lack of evidence for association rather than proof of its absence. Accordingly, we frame this study as \emph{exploratory and characterising} rather than confirmatory and avoid strong inferential claims. The released replication package \cite{our_dataset} enables future work to retest these relationships with larger samples and repeated runs.

\subsection{Future Work}

We identify several directions for future research:

\begin{enumerate}
    \item \textbf{Longitudinal study:} Track model improvements over time on fixed tasks.
    \item \textbf{Human comparison:} Measure human developer performance for baseline comparison.
    \item \textbf{Multi-language:} Extend to Java, JavaScript, Go, and other languages.
    \item \textbf{Team simulation:} Evaluate multi-agent collaboration scenarios.
    \item \textbf{Real-world validation:} Deploy models on actual PRs and measure acceptance rates.
\end{enumerate}
